\hnsection{GENESIS}

The theory, called for nearly a century ``theory of Lie groups'', was essentially developed by one mathematician: Sophus Lie.

Before embarking on the history, we summarize briefly various earlier research which prepared the way.

\textbf{(a) Transformation groups (Klein-Lie, 1869--1872)}

About 1860 the theory of permutation groups of a finite set was developing and beginning to be used (Serret, Kronecker, Mathieu, Jordan). On the other hand, the theory of invariants, which was in full flight, was familiarizing mathematicians with some infinite sets of geometric transformations stable under composition (notably linear or projective transformations). But before the work of Jordan {[}7{]} in 1868 on ``groups of movements'' (closed subgroups of the group of displacements of 3-dimensional Euclidean space), no conscious link seems to have been established between these two currents of ideas.

In 1869, the young Felix Klein (1849--1925), a pupil of Plücker, established a friendship in Berlin with the Norwegian Sophus Lie (1842--1899), who was a few years older, brought together by their common interest in Plücker's ``line geometry'' and in particular the theory of line complexes. It was about this time that Lie conceived one of his most original ideas, the introduction of the notion of invariant into Analysis and Differential Geometry; one of the sources was his observation that the classical methods of integration ``by quadratures'' of differential equations depended entirely on the fact that the equation is invariant under a ``continuous'' family of transformations. 1869 is the date of the first work (edited by Klein) where Lie used this idea; there he studied the ``Reye complex'' (set of four lines cutting the faces of a tetrahedron in four points with a given cross-ratio) and the curves and surfaces which admit as tangents lines of this complex {[}3a{]}: his method depends on the invariance of the Reye complex under the 3-parameter commutative group (maximal torus of PGL(4, C)) which leaves the vertices of the tetrahedron invariant. The same idea dominated the work Klein and Lie wrote together when they were in Paris in the spring of 1870 {[}1, a{]}; there they essentially determined the commutative connected subgroups of the projective group of the plane PGL(3, C) and studied the geometric properties of their orbits (under the name of curves or surfaces V); this gave them, by a uniform procedure, properties of various curves, algebraic or transcendental, such as \(y = cx^m\) or the logarithmic spirals. Both works served to underline the profound impression made on them by the theories of Galois and Jordan (Jordan's commentary on Galois had appeared in Math. Annalen in 1869; on the other hand, Lie had heard talk of Galois theory as early as 1863). Klein, who in 1871 became interested in non-Euclidean geometries, saw there the start of his research on a classification principle for all known geometries which was to lead him in 1872 to the ``Erlangen Programme''. For his part, Lie, in a letter of 1873 to A. Mayer ({[}3{]}, vol.~V, p.~584), dated his ideas on transformation groups from his sojourn in Paris and in a work of 1871 ({[}3 b{]} p.~208) he was already using the term ``transformation group'' and explicitly posed the problem of the determination of all subgroups (``continuous or discontinuous'') of GL(n, C). To be truthful, Klein and Lie had both experienced some difficulty in entering this new mathematical universe and Klein spoke of Jordan's newly appeared ``Treatise'' as a ``book sealed with seven seals'' ({[}2{]}, p.~51); he wrote moreover concerning {[}1 a{]} and {[}1 b{]}: ``To Lie belongs all the credit for the heuristic idea of a continuous group of operators, in particular everything concerning integration of differential equations and partial derivatives. All the notions he developed later in his theory of continuous groups were already there in embryo, but were however so little elaborated, that only after long conversations could I convince him of many details, for example to begin with the very existence of the curves V'' ({[}2{]}, p.~415).

\textbf{(b) Infinitesimal transformations}

The concept of an ``infinitely small'' transformation goes back at least to the beginnings of Infinitesimal Calculus; we know that Descartes discovered the instantaneous centre of rotation when admitting that ``in the infinitely small'' every plane movement can be likened to a rotation: the elaboration of Analytical Mechanics in the 18th century is entirely founded on similar ideas. In 1851, Sylvester, seeking to form invariants of the linear group \(\mathbf{GL}(3,\mathbf{C})\) and some of its subgroups, gave the parameters appearing in his matrices ``infinitely small'' increases of the forms \(\alpha_{j}dt\) and expressed the fact that a function \(f((z_{j}))\) was invariant by writing the equation \(f((z_{j}+\alpha_{j}(dt))=f((z_{j}))\) ; this gave him for f the linear partial differential equation Xf=0, where

\[
\mathrm{X} f = \sum_ {j} \alpha_ {j} \frac {\partial f}{\partial z _ {j}},\tag{1}
\]

where X is therefore a differential operator, the ``derivative in the direction of the direction parameters \(\alpha_{j}\) ({[}5{]}, vol.~3, pp.~326 and 327); Sylvester seemed to think that here was a general principle of considerable importance but appears never to have returned to the question. A little later, Cayley ({[}6{]}, vol.~II, pp.~164--178) proceeded similarly for the invariants of \(\mathbf{SL}(2,\mathbf{C})\) under certain representations of this group and showed that they are the solutions of two first order partial differential equations \(Xf=0\) , \(Yf=0\) , where X and Y are obtained as above from ``infinitely small'' transformations

\[
\left( \begin{array}{c c} 0 & 0 \\ d t & 0 \end{array} \right) \quad \text { and } \quad \left( \begin{array}{c c} 0 & d t \\ 0 & 0 \end{array} \right).
\]

In modern terms, this is expressed by the fact that X and Y generate the Lie algebra \(\mathfrak{sl}(2,\mathbf{C})\) ; moreover Cayley calculated the bracket XY -- YX explicitly and showed that it was also derived from an ``infinitely small'' transformation.

In his memoir of 1868 on groups of movements {[}7{]}, Jordan used from beginning to end the concept of ``infinitely small transformation'', but exclusively from a geometric point of view. He is no doubt responsible for the idea of a one-parameter group ``generated'' by an infinitely small transformation: for Jordan, it was the set of transformations obtained by ``suitably repeating'' the infinitely small transformation (loc. cit., p.~243). Klein and Lie, in their memoir, used the same expression ``repeated infinitely small transformation'' {[}1 b{]}, but the context shows that they understood by that an integration of a differential system. If the one-parameter group they considered consists of transformations \(x' = f(x, y, t)\) , \(y' = g(x, y, t)\) , the corresponding ``infinitely small transformation'' is given by

\[
d x = p (x, y) d t, \quad d y = q (x, y) d t
\]

where \(p(x,y) = \frac{\partial f}{\partial t} (x,y,t_0), q(x,y) = \frac{\partial g}{\partial t} (x,y,t_0)\) and \(t_0\) corresponds to the

identity transformation of the group. As Klein and Lie knew the functions f and g explicitly, they had no difficulty in verifying that the functions

\[
t \mapsto f (x, y, t) \quad \text { and } \quad t \mapsto g (x, y, t)
\]

give in parametric form the integral curve of the differential equation

\[
q (\xi , \eta) d \xi = p (\xi , \eta) d \eta
\]

passing through the point \((x,y)\) , but they give no general argument; moreover they nowhere use this fact in the remainder of their memoir.

\textbf{(c) Contact transformations}

In the next two years, Lie seemed to abandon the theory of transformation groups (although he remained in very close contact with Klein, who published his ``Programme'' in 1872) to study contact transformations, integration of first order partial differential equations and the relations between these two theories. We are not concerned here with the history of these questions and we shall confine ourselves to mentioning a few points which seem to have played an important role in the genesis of the theory of transformation groups.

The notion of contact transformation generalizes both point transformations and inverse polar transformations. Roughly a contact transformation† on \(\mathbf{C}^{\mathrm{n}}\) is an isomorphism of an open subset \(\Omega\) of the manifold \(T^{\prime}(\mathbf{C}^{\mathrm{n}})\) of cotangent vectors to \(\mathbf{C}^{\mathrm{n}}\) onto another open subset \(\Omega^{\prime}\) of \(T^{\prime}(\mathbf{C}^{\mathrm{n}})\) mapping the canonical 1-form of \(\Omega\) to that of \(\Omega^{\prime}\) . In other words, if \((x_{1},\ldots ,x_{n},p_{1},\ldots ,p_{n})\) denote the canonical coordinates of \(T^{\prime}(\mathbf{C}^{\mathrm{n}})\) , a contact transformation is an isomorphism \((x_{1},p_{1})\mapsto (\mathbf{X}_{\mathrm{i}},\mathbf{P}_{\mathrm{i}})\) satisfying the relation \(\sum_{i = 1}^{n}\mathbf{P}_{i}d\mathbf{X}_{i} = \sum_{i = 1}^{n}p_{i}dx_{i}\) . Such transformations occur in the study of the integration of partial differential equations of the form

\[
\mathrm{F} \left(x _ {1}, x _ {2}, \dots , x _ {n}, \frac {\partial z}{\partial x _ {1}}, \dots , \frac {\partial z}{\partial x _ {n}}\right) = 0.\tag{2}
\]

Lie became familiar during his research on these questions with the manipulation of Poisson brackets

\[
(f, g) = \sum_ {i = 1} ^ {n} \left(\frac {\partial f}{\partial x _ {i}} \frac {\partial g}{\partial p _ {i}} - \frac {\partial g}{\partial x _ {i}} \frac {\partial f}{\partial p _ {i}}\right)\tag{3}
\]

and the brackets \(\ddagger_{\xi}\) {[}X, Y{]} = XY - YX of differential operators of type (1); he interprets the Poisson bracket (3) as the effect on f of a transformation of type (1) associated with g and observes on this occasion that the Jacobi identity for Poisson brackets means that the bracket of the differential operators corresponding to f and g is associated with the bracket (g, h). Research into functions g such that (F, g) = 0, which occurs in Jacobi's method for integrating the partial differential equation (2), became for Lie the study of infinitesimal contact transformations which leave the given equation invariant. Finally, Lie was led to study sets of functions \((u_{j})_{1 \in j \leq m}\) of the \(x_{i}\) and \(p_{i}\) such that the brackets \((u_{j}, u_{k})\) are functions of the \(u_{h}\) and called these sets ``groups'' (they had already essentially been studied by Jacobi) {[}3 c{]}.

